\bta{2020}

\begin{enumerate}
	%\renewcommand{\labelenumi}{\arabic{enumi}.}
	% A(\Alph) a(\alph) I(\Roman) i(\roman) 1(\arabic)
	%设定全局标号series=example	%引用全局变量resume=example
	%[topsep=-0.3em,parsep=-0.3em,itemsep=-0.3em,partopsep=-0.3em]
	%可使用leftmargin调整列表环境左边的空白长度 [leftmargin=0em]
	\item
	选择题(本题满分 50 分，每小题 5 分)
	\begin{enumerate}
		%\renewcommand{\labelenumi}{\arabic{enumi}.}
		% A(\Alph) a(\alph) I(\Roman) i(\roman) 1(\arabic)
		%设定全局标号series=example	%引用全局变量resume=example
		%[topsep=-0.3em,parsep=-0.3em,itemsep=-0.3em,partopsep=-0.3em]
		%可使用leftmargin调整列表环境左边的空白长度 [leftmargin=0em]
		\item
		极限 $\lim\limits _{n \rightarrow+\infty}\left(\frac{2}{n^{2}}+\frac{4}{n^{2}+1}+\cdots+\frac{2 n}{n^{2}+n-1}\right)$ 的值为 \xzanswer{B} 
		
		
\fourchoices
{$ 0 $}
{$ 1 $}
{$ 2 $}
{$+\infty$}

\item 
设 $f(x)$ 在 $x=a$ 处连续, $F(x)=f(x)|x-a|$, 则 $f(a)=0$ 是 $F(x)$ 在 $x=a$ 处可导的 \xzanswer{A} 


\fourchoices
{充要条件}
{充分非必要条件}
{必要非充分条件}
{既非充分又非必要条件}

		
\item 
极限$\lim\limits _{x \rightarrow \frac{\pi}{2}}(\sin x)^{\frac{1}{1-\sin x}}$的值为 \xzanswer{A}



\fourchoices
{$e^{-1}$}
{$ 1 $}
{$e$}
{$+\infty$}

\item 
设周期函数 $f(x)$ 在 $(-\infty,+\infty)$ 内可导, 周期为 4, 又 $\lim\limits _{x \rightarrow 0} \frac{f(1)-f(1-x)}{2 x}=-1$, 则曲线 $y=f(x)$ 在点 $(5, f(5))$ 处的切线斜率为 \xzanswer{D} 


\fourchoices
{$\frac{1}{2}$}
{$ 2 $}
{$-1$}
{$-2$}

\item 
设向量 $\vec{a}=(1,2,2),  \vec{b}=(0,1,2)$, 则向量 $\vec{b}$ 在向量 $\vec{a}$ 方向上的投影向量为 \xzanswer{D} 


\fourchoices
{$\left(0, \frac{6}{5}, \frac{12}{5}\right)$}
{$\left(\frac{1}{3}, \frac{2}{3}, \frac{2}{3}\right)$}
{$\left(0, \frac{1}{\sqrt{5}}, \frac{2}{\sqrt{5}}\right)$}
{$\left(\frac{2}{3}, \frac{4}{3}, \frac{4}{3}\right)$}

\item 
二元函数 $f(x, y)$ 在点 $(0,0)$ 处可微的一个充分条件是 \xzanswer{C} 


\fourchoices
{$\lim\limits _{(x, y) \rightarrow(0,0)}[f(x, y)-f(0,0)]=0$.}
{$\lim\limits _{x \rightarrow 0} \frac{f(x, 0)-f(0,0)}{x}=0$, 且 $\lim\limits _{y \rightarrow 0} \frac{f(0, y)-f(0,0)}{y}=0$.}
{$\lim\limits _{(x, y) \rightarrow(0,0)} \frac{f(x, y)-f(0,0)}{\sqrt{x^{2}+y^{2}}}=0$}
{$\lim\limits _{x \rightarrow 0}\left[f_{x}^{\prime}(x, 0)-f_{x}^{\prime}(0,0)\right]=0$, 且 $\lim\limits _{y \rightarrow 0}\left[f_{y}^{\prime}(0, y)-f_{y}^{\prime}(0,0)\right]=0$.}
		
\item 
级数 $\sum\limits_{n=1}^{\infty} \frac{1}{n(n+1)(n+2)}=$  \xzanswer{B} 


\fourchoices
{$ 0 $}
{$\frac{1}{4}$}
{$\frac{1}{3}$}
{$ 1 $}

\item 
设方程 $x+y+z=e^{x y}$, 则 $\frac{\partial^{2} z}{\partial x^{2}}+\frac{\partial^{2} z}{\partial y^{2}}$ 的表达式是 \xzanswer{A} 


\fourchoices
{$\left(x^{2}+y^{2}\right) e^{x y}$}
{$(x+y) e^{x y}$}
{$2 x y e^{x y}$}
{$(1+x y) e^{x y}$}

\item 
设 $a_{0}=3, a_{1}=5$, 且对任何自然数 $n>1$ 有 $n a_{n}=\frac{2}{3} a_{n-1}-(n-1) a_{n-1}$, 则幂级数 $\sum\limits_{n=0}^{\infty} a_{n} x^{n}$ 的收敛半径为 \xzanswer{B} 


\fourchoices
{$\frac{2}{3}$}
{$ 1 $}
{$\frac{3}{2}$}
{$ 2 $}

\item 
下列反常积分发散的是 \xzanswer{C} 


\fourchoices
{$\int_{0}^{+\infty} \frac{x^{6}}{1+e^{x}} d x$}
{$\int_{0}^{+\infty} \frac{1}{\sqrt{x}(1+x)} d x$}
{$\int_{1}^{+\infty} \frac{1}{x^{4} \ln x} d x$}
{$\int_{1}^{+\infty} \frac{2 \sin ^{2} x}{1+x^{2}} d x$}

		
	\end{enumerate}


\banswer{
\begin{enumerate}
	%\renewcommand{\labelenumi}{\arabic{enumi}.}
	% A(\Alph) a(\alph) I(\Roman) i(\roman) 1(\arabic)
	\item
	\textbf{B}。记 $S_n=\sum\limits_{k=1}^{n}\frac{2k}{n^2+k-1}$。由 $n^2\leqslant n^2+k-1\leqslant n^2+n-1$
	及 $\sum\limits_{k=1}^{n}2k=n(n+1)$ 得
	\[
	\frac{n(n+1)}{n^2+n-1}\leqslant S_n\leqslant\frac{n(n+1)}{n^2} .
	\]
	两端当 $n\to\infty$ 时均趋于 $1$，由夹逼准则 $S_n\to1$。

	\item
	\textbf{A}。因 $F(a)=f(a)\cdot0=0$，故
	\[
	\frac{F(x)-F(a)}{x-a}=f(x)\cdot\frac{|x-a|}{x-a} .
	\]
	由 $f$ 在 $x=a$ 连续，上式右极限为 $f(a)$、左极限为 $-f(a)$；
	二者相等当且仅当 $f(a)=0$。故 $f(a)=0$ 是 $F$ 在 $x=a$ 可导的充要条件。

	\item
	\textbf{A}。因 $\sin x\leqslant1$，令 $u=1-\sin x$，则 $x\to\frac{\pi}{2}$ 时 $u\to0^{+}$，于是
	\[
	(\sin x)^{\frac{1}{1-\sin x}}=(1-u)^{\frac1u}\longrightarrow e^{-1} .
	\]

	\item
	\textbf{D}。由题设 $\lim\limits_{x\to0}\frac{f(1)-f(1-x)}{x}=-2$，
	而该式即 $\lim\limits_{x\to0}\frac{f(1-x)-f(1)}{-x}=f^{\prime}(1)$，故 $f^{\prime}(1)=-2$。
	$f$ 以 $4$ 为周期，对 $f(x+4)=f(x)$ 两端求导知 $f^{\prime}$ 也以 $4$ 为周期，
	故所求斜率 $f^{\prime}(5)=f^{\prime}(1)=-2$。

	\item
	\textbf{D}。投影向量为
	\[
	\frac{\vec{a}\cdot\vec{b}}{|\vec{a}|^{2}}\,\vec{a}
	=\frac{0+2+4}{9}(1,2,2)=\left(\frac23,\frac43,\frac43\right).
	\]
	（选项 C 只是 $\vec{a}$ 方向上的单位向量之误，A 把投影方向取成了 $\vec{b}$。）

	\item
	\textbf{C}。C 即 $f(x,y)-f(0,0)=o\left(\sqrt{x^2+y^2}\right)$，
	按定义 $f$ 在 $(0,0)$ 可微（且全微分为 $0$），故充分。
	A 只给出连续性；B 只说明两个偏导数为 $0$（例如 $f=\sqrt{|xy|}$ 满足 B 却不可微）；
	D 只涉及偏导数沿两条坐标轴方向的极限，也推不出可微。

	\item
	\textbf{B}。裂项
	\[
	\frac{1}{n(n+1)(n+2)}=\frac12\left[\frac{1}{n(n+1)}-\frac{1}{(n+1)(n+2)}\right],
	\]
	部分和 $S_N=\frac12\left[\frac12-\frac{1}{(N+1)(N+2)}\right]\to\frac14$。

	\item
	\textbf{A}。由 $x+y+z=e^{xy}$ 直接解出 $z=e^{xy}-x-y$，于是
	$z_x=ye^{xy}-1$，$z_{xx}=y^2e^{xy}$；由对称性 $z_{yy}=x^2e^{xy}$。故
	$z_{xx}+z_{yy}=\left(x^2+y^2\right)e^{xy}$。

	\item
	\textbf{B}。递推式 $na_n=\left(\frac23-n+1\right)a_{n-1}$ 给出
	$a_n=\frac{\frac53-n}{n}a_{n-1}$，故
	\[
	\left|\frac{a_n}{a_{n-1}}\right|=\frac{\left|n-\frac53\right|}{n}\longrightarrow1 ,
	\]
	由比值法收敛半径 $R=1$。

	\item
	\textbf{C}。C 中 $x\to1^{+}$ 时 $\ln x\sim x-1$，故 $\frac{1}{x^4\ln x}\sim\frac{1}{x-1}$，
	瑕积分发散。
	A 中 $\frac{x^6}{1+e^x}\leqslant x^6e^{-x}$，收敛；
	B 中令 $x=t^2$ 得 $\int_0^{+\infty}\frac{2\,dt}{1+t^2}=\pi$，收敛；
	D 中 $\frac{2\sin^2x}{1+x^2}\leqslant\frac{2}{1+x^2}$，收敛。
\end{enumerate}

	\par\vspace{1mm}\noindent
	\textbf{注：}{\kaishu
	第 (3) 题原答案作 D（$+\infty$），有误。底数 $\sin x\to1^{-}$、指数 $\to+\infty$，
	是 $1^{\infty}$ 型而非 $+\infty$ 型；取 $x=\frac{\pi}{2}-10^{-4}$ 数值计算得 $0.367879$，
	正是 $e^{-1}$。正确答案为 A。
	}\par
}


\item
（本题满分10分）
已知 $\lim\limits _{x \rightarrow 0} f(x)$ 存在,而 $\lim\limits _{x \rightarrow 0} f^{\prime}(x)$ 不存在,并且 $\lim\limits _{x \rightarrow 0} \frac{\sqrt[3]{1+x f(x)}-1}{\sin x}=3$，求 $\lim\limits _{x \rightarrow 0} f(x)$。

\banswer{
	设 $\lim\limits_{x\to0}f(x)=A$（题设该极限存在），则 $xf(x)\to0$。
	由等价无穷小 $\sqrt[3]{1+u}-1\sim\frac{u}{3}\ (u\to0)$ 与 $\sin x\sim x$ 得
	\[
	\frac{\sqrt[3]{1+xf(x)}-1}{\sin x}\sim\frac{\frac13\,xf(x)}{x}=\frac{f(x)}{3}
	\longrightarrow\frac{A}{3} .
	\]
	由题设该极限为 $3$，故 $\frac{A}{3}=3$，即
	\[
	\lim_{x\to0}f(x)=9 .
	\]

	\par\vspace{1mm}\noindent
	\textbf{注：}{\kaishu
	题中“$\lim\limits_{x\to0}f^{\prime}(x)$ 不存在”只是提醒不能用洛必达法则，
	上面的等价无穷小替换并不需要 $f$ 可导。
	}\par
}


\item 
（本题满分10分）
两平面均通过点 $A(-2,1,-1)$, 其中一个平面通过 $x$ 轴, 另一个平面通过直线 $\frac{x-1}{2}=\frac{y+1}{1}=\frac{z}{-1}$ ，求两平面夹角的余弦。

\banswer{
	\textbf{定过 $x$ 轴的平面 $\pi_1$.} 过 $x$ 轴的平面不含 $x$ 项且过原点，方程形如 $By+Cz=0$。
	代入 $A(-2,1,-1)$ 得 $B-C=0$，取 $B=C=1$，即
	\[
	\pi_1:\ y+z=0,\qquad \bm{n}_1=(0,1,1).
	\]

	\textbf{定过已知直线的平面 $\pi_2$.} 该直线过点 $M(1,-1,0)$，方向向量 $\bm{s}=(2,1,-1)$；
	又 $\overrightarrow{MA}=(-3,2,-1)$，故 $\pi_2$ 的法向量
	\[
	\bm{n}_2=\bm{s}\times\overrightarrow{MA}
	=\begin{vmatrix}\bm{i}&\bm{j}&\bm{k}\\2&1&-1\\-3&2&-1\end{vmatrix}
	=(1,5,7).
	\]

	\textbf{求夹角.} $\bm{n}_1\cdot\bm{n}_2=0+5+7=12$，$\left|\bm{n}_1\right|=\sqrt2$，
	$\left|\bm{n}_2\right|=\sqrt{1+25+49}=5\sqrt3$，故
	\[
	\cos\theta=\frac{\left|\bm{n}_1\cdot\bm{n}_2\right|}{\left|\bm{n}_1\right|\left|\bm{n}_2\right|}
	=\frac{12}{5\sqrt6}=\frac{2\sqrt6}{5}\approx0.9798 .
	\]
}



\item 
（本题满分10分）
设函数 $y=f(x)$ 由 $\left\{\begin{array}{c}x^{x}+t x-t^{2}=0, \\ \arctan (t y)=\ln \left(1+t^{2} y^{2}\right)\end{array}\right.$确定，求 $\frac{d y}{d x}$ 。

\banswer{
	方程组以 $t$ 为参数确定 $x=x(t)$、$y=y(t)$，故 $\frac{dy}{dx}=\dfrac{dy/dt}{dx/dt}$。

	\textbf{由第一式定 $\frac{dx}{dt}$.} 注意 $\left(x^{x}\right)^{\prime}_x=x^{x}(1+\ln x)$，
	对 $x^{x}+tx-t^{2}=0$ 关于 $t$ 求导：
	\[
	\left[x^{x}(1+\ln x)+t\right]\frac{dx}{dt}+x-2t=0
	\;\Longrightarrow\;
	\frac{dx}{dt}=\frac{2t-x}{x^{x}(1+\ln x)+t} .
	\]

	\textbf{由第二式定 $\frac{dy}{dt}$.} 记 $w=ty$，第二式即 $\arctan w=\ln\left(1+w^{2}\right)$，
	两端对 $t$ 求导得
	\[
	\frac{w^{\prime}}{1+w^{2}}=\frac{2w\,w^{\prime}}{1+w^{2}}
	\;\Longrightarrow\;
	w^{\prime}(1-2w)=0 .
	\]
	若 $1-2w\neq0$ 则 $w^{\prime}=0$；若 $1-2w=0$ 则 $w\equiv\frac12$ 为常数，仍有 $w^{\prime}=0$。
	总之 $(ty)^{\prime}=y+t\frac{dy}{dt}=0$，即 $\frac{dy}{dt}=-\frac{y}{t}$。

	\textbf{合成.}
	\[
	\frac{dy}{dx}=-\frac{y}{t}\cdot\frac{x^{x}(1+\ln x)+t}{2t-x}
	=\frac{y\left[t+x^{x}(1+\ln x)\right]}{t(x-2t)} .
	\]

	\par\vspace{1mm}\noindent
	\textbf{注：}{\kaishu
	原答案作 $\frac{t+x^{x}(1+\ln x)}{2t^{2}(x-2t)}$，相当于在上式中取 $y=\frac{1}{2t}$，
	即把推导中必须讨论的分支 $1-2ty=0$ 当成了解，有误：
	$\arctan\frac12=0.4636\neq\ln\frac54=0.2231$，$ty=\frac12$ 并不满足第二式。
	另外由 $(ty)^{\prime}=0$ 知 $ty\equiv C$ 为常数，$C$ 由 $\arctan C=\ln\left(1+C^{2}\right)$ 定出，
	只有 $C=0$（此时 $y\equiv0$，$\frac{dy}{dx}=0$）与 $C\approx1.1720$ 两支；
	若要写成只含 $x$、$t$ 的式子，应为
	$\frac{dy}{dx}=\frac{C\left[t+x^{x}(1+\ln x)\right]}{t^{2}(x-2t)}$。
	}\par
}



\item 
（本题满分10分）
已 知 函 数 $u=f(r), r=\ln \sqrt{x^{2}+y^{2}+z^{2}}$ 满 足 方 程$\frac{\partial^{2} u}{\partial x^{2}}+\frac{\partial^{2} u}{\partial y^{2}}+\frac{\partial^{2} u}{\partial z^{2}}=\left(x^{2}+y^{2}+z^{2}\right)^{-3 / 2}$, 求 $f(x)$ 的表达式。


\banswer{
	记 $\rho=\sqrt{x^{2}+y^{2}+z^{2}}$，则 $r=\ln\rho$、$\rho=e^{r}$，$u$ 是球对称函数。

	\textbf{化简 Laplace 算子.} 对球对称函数
	$\Delta u=\frac{1}{\rho^{2}}\frac{\partial}{\partial\rho}\left(\rho^{2}\frac{\partial u}{\partial\rho}\right)$。
	由 $\frac{\partial u}{\partial\rho}=f^{\prime}(r)\frac{dr}{d\rho}=\frac{f^{\prime}(r)}{\rho}$
	得 $\rho^{2}u_\rho=\rho f^{\prime}(r)$，再对 $\rho$ 求导得
	$f^{\prime}(r)+\rho f^{\prime\prime}(r)\frac{1}{\rho}=f^{\prime}(r)+f^{\prime\prime}(r)$，故
	\[
	\Delta u=\frac{f^{\prime\prime}(r)+f^{\prime}(r)}{\rho^{2}}
	=\frac{f^{\prime\prime}(r)+f^{\prime}(r)}{e^{2r}} .
	\]

	\textbf{解常微分方程.} 题设右端 $\left(x^{2}+y^{2}+z^{2}\right)^{-3/2}=\rho^{-3}=e^{-3r}$，故
	\[
	f^{\prime\prime}(r)+f^{\prime}(r)=e^{-r} .
	\]
	特征方程 $\lambda^{2}+\lambda=0$ 的根为 $0$ 与 $-1$，齐次解为 $C_1+C_2e^{-r}$。
	因 $-1$ 是单根，设特解 $f^{*}=Are^{-r}$，代入左端得 $-Ae^{-r}$，故 $A=-1$。于是
	\[
	f(r)=C_1+C_2e^{-r}-re^{-r},
	\]
	把自变量改记为 $x$ 即 $f(x)=C_1+\left(C_2-x\right)e^{-x}$。

	\textbf{回代检验.} 此时 $u=C_1+\frac{C_2}{\rho}-\frac{\ln\rho}{\rho}$。
	直接计算 $\Delta\left(\frac1\rho\right)=0$，而
	\[
	\Delta\left(-\frac{\ln\rho}{\rho}\right)
	=g^{\prime\prime}(\rho)+\frac{2}{\rho}g^{\prime}(\rho)=\frac{1}{\rho^{3}},
	\qquad g(\rho)=-\frac{\ln\rho}{\rho},
	\]
	与题设方程右端一致。

	\par\vspace{1mm}\noindent
	\textbf{注：}{\kaishu
	原答案作 $f(x)=1+C_2e^{x}+\frac12\frac{1}{x^2+y^2+z^2}-\frac{1}{\sqrt{x^2+y^2+z^2}}$，
	有误：$f$ 是一元函数，右端不应出现 $x,y,z$；指数应为 $e^{-x}$；
	而 $\frac{1}{2\rho^{2}}$ 一项并不调和，也不是所需的特解。
	}\par
}


\item 
（本题满分10分）
设曲线 $C: y=x^{3}+2 x$ 与其在（1,3）点处的切线以及 $x$ 轴围成的区 域落在第一象限中的部分为 $D$, 计算:
\begin{enumerate}
	%\renewcommand{\labelenumi}{\arabic{enumi}.}
	% A(\Alph) a(\alph) I(\Roman) i(\roman) 1(\arabic)
	%设定全局标号series=example	%引用全局变量resume=example
	%[topsep=-0.3em,parsep=-0.3em,itemsep=-0.3em,partopsep=-0.3em]
	%可使用leftmargin调整列表环境左边的空白长度 [leftmargin=0em]
	\item
$  D  $的面积。
\item 
$D$ 绕 $x$ 轴旋转一周所得旋转体的体积。
\end{enumerate}

\banswer{
	\begin{enumerate}
		%\renewcommand{\labelenumi}{\arabic{enumi}.}
		% A(\Alph) a(\alph) I(\Roman) i(\roman) 1(\arabic)
		%设定全局标号series=example	%引用全局变量resume=example
		%[topsep=-0.3em,parsep=-0.3em,itemsep=-0.3em,partopsep=-0.3em]
		%可使用leftmargin调整列表环境左边的空白长度 [leftmargin=0em]
		\item
		\textbf{求切线并定出 $D$.} $y^{\prime}=3x^{2}+2$，在 $(1,3)$ 处斜率为 $5$，切线为
		$y=3+5(x-1)=5x-2$，它与 $x$ 轴交于 $\left(\frac25,0\right)$。
		由 $x^{3}+2x=5x-2\iff(x-1)^{2}(x+2)=0$ 知切线与曲线除切点 $(1,3)$ 外只在 $x=-2$ 处相交，
		故第一象限内的区域 $D$ 是以 $(0,0)$、$\left(\frac25,0\right)$、$(1,3)$ 为顶点的曲边三角形：
		上边界为曲线，右下边界为切线，下边界为 $x$ 轴。

		\textbf{求面积（曲线下面积减去三角形面积）.}
		\begin{align*}
		S_D&=\int_0^1\left(x^{3}+2x\right)dx-\frac12\left(1-\frac25\right)\cdot3\\
		&=\left(\frac14+1\right)-\frac{9}{10}=\frac54-\frac{9}{10}=\frac{7}{20}=0.35 .
		\end{align*}

		\item
		绕 $x$ 轴旋转用垫圈法：$0\leqslant x\leqslant\frac25$ 时上、下边界为 $x^{3}+2x$ 与 $0$；
		$\frac25\leqslant x\leqslant1$ 时为 $x^{3}+2x$ 与 $5x-2$。二者可合写为
		\[
		V_D=\pi\int_0^1\left(x^{3}+2x\right)^{2}dx-\pi\int_{2/5}^{1}(5x-2)^{2}dx .
		\]
		其中
		\begin{align*}
		\int_0^1\left(x^{6}+4x^{4}+4x^{2}\right)dx&=\frac17+\frac45+\frac43=\frac{239}{105},\\
		\int_{2/5}^{1}(5x-2)^{2}dx&=\left[\frac{(5x-2)^{3}}{15}\right]_{2/5}^{1}=\frac{27}{15}=\frac95 .
		\end{align*}
		故
		\[
		V_D=\pi\left(\frac{239}{105}-\frac{189}{105}\right)=\frac{10}{21}\pi\approx0.48\pi .
		\]

	\end{enumerate}

	\par\vspace{1mm}\noindent
	\textbf{注：}{\kaishu
	原答案作 $S_D=\frac{3}{20}$、$V_D=0.11\pi$，均有误；
	正确值为 $S_D=\frac{7}{20}$、$V_D=\frac{10}{21}\pi\approx0.48\pi$，数值积分亦予证实。
	}\par
}



\newpage
\item 
（本题满分10分）
计算下列第二型曲面积分:
\[
I=\iint_{S} x d y d z+2 y^{4} d x d z+3 z^{6} d x d y
\]
其中 $S$ 是椭球面： $x^{2}+4 y^{2}+9 z^{2}=1$ 。

\banswer{
	$S$ 是闭曲面，取外侧（题面未注明侧，按惯例取外侧）；记 $\Omega$ 为椭球体
	$x^{2}+4y^{2}+9z^{2}\leqslant1$。注意 $dxdz=-dzdx$，先把积分化为标准记法
	\[
	I=\iint_S x\,dydz-2y^{4}\,dzdx+3z^{6}\,dxdy .
	\]
	由高斯公式
	\begin{align*}
	I&=\iiint_\Omega\left[\frac{\partial x}{\partial x}
	-\frac{\partial\left(2y^{4}\right)}{\partial y}
	+\frac{\partial\left(3z^{6}\right)}{\partial z}\right]dV\\
	&=\iiint_\Omega\left(1-8y^{3}+18z^{5}\right)dV .
	\end{align*}
	$\Omega$ 关于 $zOx$ 面对称而 $y^{3}$ 关于 $y$ 为奇函数，故 $\iiint_\Omega y^{3}dV=0$；
	同理 $\iiint_\Omega z^{5}dV=0$。于是 $I$ 就是 $\Omega$ 的体积：
	\[
	I=\frac43\pi abc=\frac43\pi\cdot1\cdot\frac12\cdot\frac13=\frac{2}{9}\pi .
	\]

	\par\vspace{1mm}\noindent
	\textbf{注：}{\kaishu
	由于 $8y^{3}$、$18z^{5}$ 的积分都为零，把题面的 $dxdz$ 读作 $dzdx$ 并不改变结果，
	仍为 $\frac29\pi$。
	}\par
}




\item 
（本题满分10分）
设 $f(x)$ 是周期为$  3  $的连续函数, 证明：在任意长度为$  2  $的闭区间
$[a, a+2]$ 上至少存在一点 $\theta$, 使得 $f(\theta)=f(\theta+1)$ 。


\banswer{
	令 $g(x)=f(x)-f(x+1)$。由 $f$ 连续知 $g$ 在 $\mathbb{R}$ 上连续。
	\begin{align*}
	&g(a)+g(a+1)+g(a+2)\\
	&\quad=[f(a)-f(a+1)]+[f(a+1)-f(a+2)]+[f(a+2)-f(a+3)]\\
	&\quad=f(a)-f(a+3).
	\end{align*}
	因 $f$ 以 $3$ 为周期，$f(a+3)=f(a)$，故 $g(a)+g(a+1)+g(a+2)=0$。
	若 $g(a),g(a+1),g(a+2)$ 均为零，则任取其一为 $\theta$ 即得结论。
	否则三者不能同号（否则和不为零），故存在 $p,q\in\{a,a+1,a+2\}$ 使 $g(p)g(q)\leq0$。
	由连续函数的介值定理，$g$ 在 $p$ 与 $q$ 之间的区间上有零点，该零点落在 $[a,a+2]$ 内，记为 $\theta$，即有 $f(\theta)=f(\theta+1)$。
}


\item 
(本题满分 10 分）
设 $f(x), g(x)$ 在 $[a, b]$ 上二阶可导, 且 $f(a)=f(b)=g(a)=0$. 证明: 存在 $\xi \in(a, b)$, 使得 $f^{\prime \prime}(\xi) g(\xi)+2 f^{\prime}(\xi) g^{\prime}(\xi)+f(\xi) g^{\prime \prime}(\xi)=0$ 。


\banswer{
	令 $h(x)=f(x)g(x)$，则 $h(a)=f(a)g(a)=0$，$h(b)=f(b)g(b)=0$。
	由 Rolle 定理，存在 $\eta\in(a,b)$ 使 $h'(\eta)=0$。
	又 $h'(x)=f'(x)g(x)+f(x)g'(x)$，故 $h'(a)=f'(a)g(a)+f(a)g'(a)=0$。
	于是 $h'(a)=h'(\eta)=0$ 且 $a<\eta$，对 $h'$ 在 $[a,\eta]$ 上再用 Rolle 定理，
	存在 $\xi\in(a,\eta)\subset(a,b)$ 使 $h''(\xi)=0$。
	而 $h''(x)=f''(x)g(x)+2f'(x)g'(x)+f(x)g''(x)$，代入 $\xi$ 即得所证。
}


\item 
（本题满分10分）
设 $f$ 是 $[0,1]$ 上的连续函数，满足 $\int_{0}^{x} f(t) d t \geq 0$ 对所有的 $x \in[0,1]$ 成立且 $\int_{0}^{1} f(t) d t=0$ 。证明: $\int_{0}^{1} x f(x) d x \leq 0$ 。


\banswer{
	令 $F(x)=\int_0^x f(t)\,dt$，则 $F(0)=0$，$F(1)=0$，$F(x)\geq0\;(x\in[0,1])$，且 $F'(x)=f(x)$。
	由分部积分，
	\begin{align*}
	\int_0^1 x f(x)\,dx&=\int_0^1 x\,F'(x)\,dx
	=\bigl[xF(x)\bigr]_0^1-\int_0^1 F(x)\,dx\\
	&=F(1)-\int_0^1 F(x)\,dx=-\int_0^1 F(x)\,dx.
	\end{align*}
	因 $F(x)\geq0$，故 $\int_0^1 F(x)\,dx\geq0$，从而 $\int_0^1 x f(x)\,dx\leq0$。
}



\item 
（本题满分10分）
求证：若正数 $x, y, z$ 满足 $x^{2}+y^{2}+z^{2}=a$, 其中 $a>0$, 则有不等
式 $x^{3}+y^{3}+z^{3} \geq \frac{a \sqrt{3 a}}{3}$ 恒成立。

\banswer{
	由 Cauchy–Schwarz 不等式，
	\[
	(x^3+y^3+z^3)(x+y+z)\geq(x^2+y^2+z^2)^2=a^2,
	\]
	故 $x^3+y^3+z^3\geq\dfrac{a^2}{x+y+z}$。
	再由 Cauchy–Schwarz，$(x+y+z)^2\leq3(x^2+y^2+z^2)=3a$，得 $x+y+z\leq\sqrt{3a}$，
	从而 $\dfrac{1}{x+y+z}\geq\dfrac{1}{\sqrt{3a}}$。于是
	\[
	x^3+y^3+z^3\geq\frac{a^2}{\sqrt{3a}}=\frac{a\sqrt{3a}}{3}.
	\]
}

	
\end{enumerate}


